\section{The source paper and replication scope}
\label{sec:source}

The model replicated here is documented in \citet{brignone2025}: Davide Brignone and Michele Piffer, \emph{A structural VAR model for the UK economy}, Bank of England Macro Technical Paper No.~3, published 18 July 2025. It is one of the first papers in the Bank's Macro Technical Paper series, which documents the modelling toolkit supporting the Monetary Policy Committee, and it describes the SVAR used on a round-by-round basis to provide a structural narrative for revisions to the Monetary Policy Report forecast. A companion paper by the same authors --- Staff Working Paper No.~1{,}165, January 2026 --- develops the forecast-revision methodology in general form \citep{brignone2026}; Section~\ref{sec:comparison-swp} compares the replication against it.

\paragraph{Replication scope.} The replication targets the paper's core pipeline and its standard outputs (the paper's Figures 2--6): Bayesian estimation, identification, impulse responses, FEVDs, estimated structural shocks and historical decompositions, plus the Section-5 forecast-revision machinery in reduced scope. The estimation sample matches the paper exactly: \textbf{1992Q1--2023Q2} (126 quarters), beginning with UK inflation targeting and ending a year before the data edge because recent quarters are revision-prone. Data through \textbf{2024Q2} are retained for the forecasting and forecast-revision exercises, matching the vintage of the Bank's August 2024 Monetary Policy Report round \citep{boe2024mpr}.

\paragraph{Approximated internal series.} Two of the eight variables --- UK-trade-weighted world GDP and world CPI --- are internal Bank series that are not published. The replication proxies them with trade-weighted aggregates constructed from OECD and IMF sources, and constructs the real sterling oil price from Brent crude converted to sterling and deflated by UK CPI, since the paper does not fully specify its construction. UK series (Bank Rate, sterling ERI, CPISA, CPI Energy, real GDP) come from the Bank of England and the Office for National Statistics \citep{ons2025gdp, ons2024cpi}. Because the world aggregates are proxies and the Bank released no replication package, the global block can be validated only against the published figures, not against the underlying data --- a caveat the quantitative comparison in Section~\ref{sec:validation} makes explicit.

\paragraph{Documented deviations.} Beyond the proxy world series, the replication documents its deviations openly: an assumed lag length of $p=4$ (unstated in the paper), default rather than hierarchically optimised prior hyperparameters, Covid dummies as exogenous regressors rather than the full pandemic-prior algebra, permutation search restricted to zero-pattern-equivalent groups, and current data vintages in place of the August 2024 vintages. Each deviation is a configurable parameter or a flagged approximation in the codebase, not a silent choice.
